\documentclass[11pt]{article}
\usepackage[margin=1.1in]{geometry}
\usepackage{amsmath,amssymb,amsthm,mathtools}
\usepackage[hidelinks]{hyperref}

\newtheorem{theorem}{Theorem}
\newtheorem{lemma}{Lemma}
\newtheorem{proposition}{Proposition}
\newtheorem{remark}{Remark}
\newcommand{\R}{\mathbb R}
\newcommand{\boxt}{\boxtimes}

\title{Failure of Free Multiplicative Strong Log-Unimodality at Every Time}
\author{}
\date{September 2026}

\begin{document}
\maketitle

\begin{abstract}
Let $\sigma_t$ denote the marginal law of free positive multiplicative Brownian motion. Hasebe, Ueda and Wang asked whether $\sigma_t$ fails free multiplicative strong log-unimodality for every $t>0$. We answer this affirmatively. For each $t>0$ we give an explicit log-unimodal probability measure $\nu_t$ on $(0,\infty)$ such that $\sigma_t\boxtimes\nu_t$ is not log-unimodal. The proof uses the root-count criterion of Hasebe--Ueda--Wang and a one-parameter family of Poisson-type kernels whose normalized profiles are monotone in the kernel scale.
\end{abstract}

\section{Statement and source criterion}

A probability measure $\mu$ on $(0,\infty)$ is \emph{log-unimodal} when $\log_*\mu$ is unimodal on $\R$. It is \emph{$\boxtimes$-strongly log-unimodal} when $\mu\boxtimes\nu$ is log-unimodal for every log-unimodal probability measure $\nu$ on $(0,\infty)$.

Hasebe, Ueda and Wang~\cite{HUW} proved that $\sigma_t$ fails this strong property for sufficiently small $t>0$ and asked in their Problem~4.11 whether the failure holds for every $t>0$. We prove the full statement.

\begin{theorem}\label{thm:main}
For every $t>0$, the free positive multiplicative Brownian law $\sigma_t$ is not $\boxtimes$-strongly log-unimodal. More precisely, for every $t>0$ there is an explicit log-unimodal probability measure $\nu_t$ such that $\sigma_t\boxtimes\nu_t$ is not log-unimodal.
\end{theorem}

We use Hasebe--Ueda--Wang's Lemma~4.1 in the following contrapositive form. For $b\in(0,\pi)$ let
\begin{equation}\label{eq:lambda}
 \lambda_b(dr)=\frac{c_b}{1-2r\cos b+r^2}\,dr,
 \qquad c_b=\frac{\sin b}{\pi-b}.
\end{equation}
If $\sigma_t\boxtimes\nu$ is log-unimodal, then for every $a\in(0,1/t)$ the equation
\begin{equation}\label{eq:source-level}
 \frac{r}{c_{a\pi t}}\frac{d(\lambda_{a\pi t}\circledast\nu^{-1})}{dr}(r)
 =\frac{a\pi}{\sin(a\pi t)}
\end{equation}
has at most two positive solutions $r$. Here $\circledast$ is classical multiplicative convolution and $\nu^{-1}$ is the image of $\nu$ under $x\mapsto x^{-1}$. Thus it suffices to construct a log-unimodal $\nu$ and one admissible $a$ for which \eqref{eq:source-level} has four positive roots.

\section{A normalized kernel}

Under logarithmic coordinates $r=e^x$, the density of $\log_*\lambda_b$ is
\begin{equation}\label{eq:h}
 h_b(x)=\frac{c_b}{2(\cosh x-\cos b)}.
\end{equation}
Set
\begin{equation}\label{eq:K}
 A_b=\frac{\pi}{\pi-b},
 \qquad
 K_b(v)=\frac{b}{A_b}h_b(bv)
 =\frac{b\sin b}{2\pi(\cosh(bv)-\cos b)}.
\end{equation}

\begin{lemma}[scale monotonicity]\label{lem:monotone}
For every fixed $v\in\R$, $b\mapsto K_b(v)$ is non-increasing on $(0,\pi)$. Moreover
\begin{equation}\label{eq:K0}
 K_b(v)\le K_0(v):=\frac{1}{\pi(1+v^2)}.
\end{equation}
\end{lemma}

\begin{proof}
Write
\[
 K_b(v)=\frac{1}{2\pi G_v(b)},
 \qquad
 G_v(b)=\frac{\cosh(bv)-\cos b}{b\sin b}.
\]
The decomposition
\begin{equation}\label{eq:G-decomp}
 G_v(b)
 =\frac{\cosh(bv)-1}{b^2}\frac{b}{\sin b}
 +\frac{\tan(b/2)}{b}
\end{equation}
shows that $G_v$ is non-decreasing. Indeed, $b/\sin b$ is increasing because
\[
 \left(\frac b{\sin b}\right)'
 =\frac{\sin b-b\cos b}{\sin^2b}>0,
\]
and $\tan(b/2)/b$ is increasing because $x\mapsto \tan x/x$ is increasing on $(0,\pi/2)$. Finally, with
\[
 \phi(x)=\frac{\cosh x-1}{x^2},
\]
we have
\[
 \phi'(x)=\frac{x\sinh x-2(\cosh x-1)}{x^3}\ge0
 \quad (x>0),
\]
because the numerator has derivative $x\cosh x-\sinh x$ and second derivative $x\sinh x\ge0$, both starting from zero at the origin. Since $\phi$ is even, $(\cosh(bv)-1)/b^2=v^2\phi(bv)$ is non-decreasing in $b>0$. Hence $G_v$ is non-decreasing and $K_b$ is non-increasing. Taking $b\downarrow0$ in \eqref{eq:K} yields \eqref{eq:K0}.
\end{proof}

An exact primitive is
\begin{equation}\label{eq:primitive}
 \int_u^v K_b(w)\,dw
 =\frac{\Phi_b(v)-\Phi_b(u)}{\pi},
 \qquad
 \Phi_b(v)=\arctan\!\left(\cot\frac b2\,\tanh\frac{bv}{2}\right).
\end{equation}

\section{The witness profile}

Fix
\begin{equation}\label{eq:constants}
 B=\frac1{10},\qquad \kappa=\frac85,\qquad d=\frac{19}{20}.
\end{equation}
For $0<b\le B$ and $\ell\ge100$, set
\[
 p=\frac{\kappa}{\ell}
\]
and define the probability measure $\nu_{b,p,\ell}$ by
\begin{equation}\label{eq:witness}
 \log_*\nu_{b,p,\ell}
 =p\delta_0+\frac{1-p}{\ell b}\mathbf 1_{[-\ell b,0]}(y)\,dy.
\end{equation}
This is unimodal on the logarithmic line with mode $0$, hence $\nu_{b,p,\ell}$ is log-unimodal with mode $1$.

The inverse law has its uniform logarithmic density on $[0,\ell b]$. Therefore the logarithmic density $q_{b,p,\ell}$ of $\lambda_b\circledast\nu_{b,p,\ell}^{-1}$ is
\[
 q_{b,p,\ell}(x)
 =p h_b(x)+\frac{1-p}{\ell b}\int_0^{\ell b}h_b(x-y)\,dy.
\]
Define
\begin{equation}\label{eq:S}
 S_{b,p,\ell}(y)
 :=\frac b{A_b}q_{b,p,\ell}(by)
 =pK_b(y)+\frac{1-p}{\ell}\int_{y-\ell}^{y}K_b(v)\,dv
\end{equation}
and
\begin{equation}\label{eq:R}
 R_{b,\ell}(y):=\ell S_{b,p,\ell}(y)
 =\kappa K_b(y)+\left(1-\frac{\kappa}{\ell}\right)
 \int_{y-\ell}^{y}K_b(v)\,dv.
\end{equation}

\begin{lemma}[uniform high--low--high pattern]\label{lem:crossings}
For every $0<b\le1/10$ and every $\ell\ge100$,
\begin{equation}\label{eq:HLH}
 R_{b,\ell}(0)>\frac{19}{20},\qquad
 R_{b,\ell}(3)<\frac{19}{20},\qquad
 R_{b,\ell}(\ell/2)>\frac{19}{20}.
\end{equation}
\end{lemma}

\begin{proof}
Since $\ell\ge100$,
\begin{equation}\label{eq:coeff}
 1-\frac{\kappa}{\ell}\ge\frac{123}{125}.
\end{equation}

At $y=0$, Lemma~\ref{lem:monotone} and interval inclusion give
\[
 R_{b,\ell}(0)
 \ge \frac85 K_B(0)+\frac{123}{125}\int_{-100}^0K_B(v)\,dv.
\]
Elementary Taylor estimates and $\pi<22/7$ give
\[
 K_B(0)>\frac{4193}{13200}.
\]
Using \eqref{eq:primitive},
\[
 \int_{-100}^0K_B(v)\,dv
 =\frac1\pi\arctan\!\left(\cot\frac1{20}\tanh5\right).
\]
The bounds $\cot(1/20)>799/40$ and $\tanh5>9/10$ make the arctangent argument larger than $16$. Since $\arctan(16)>\pi/2-1/16>12\pi/25$, we obtain
\[
 \int_{-100}^0K_B>\frac{12}{25}.
\]
Therefore
\begin{equation}\label{eq:high0}
 R_{b,\ell}(0)
 >\frac85\frac{4193}{13200}+\frac{123}{125}\frac{12}{25}
 =\frac{202241}{206250}>\frac{19}{20}.
\end{equation}

At $y=3$, Lemma~\ref{lem:monotone} gives
\[
 R_{b,\ell}(3)
 \le \frac85K_0(3)+\int_{-\infty}^3K_0(v)\,dv
 =1-\frac{\arctan(1/3)}\pi+\frac4{25\pi}.
\]
The alternating series gives $\arctan(1/3)>26/81$, while
\[
 \frac\pi{20}+\frac4{25}<\frac{11}{70}+\frac4{25}=\frac{111}{350}
\]
and
\[
 \frac{26}{81}-\frac{111}{350}=\frac{109}{28350}>0.
\]
Hence $R_{b,\ell}(3)<19/20$.

At $y=\ell/2$, discard the nonnegative first term in \eqref{eq:R}. By Lemma~\ref{lem:monotone}, \eqref{eq:coeff}, and $\ell\ge100$,
\[
 R_{b,\ell}(\ell/2)
 \ge\frac{123}{125}\int_{-50}^{50}K_B(v)\,dv.
\]
The primitive \eqref{eq:primitive} gives
\[
 \int_{-50}^{50}K_B
 =\frac2\pi\arctan\!\left(\cot\frac1{20}\tanh\frac52\right).
\]
Now $\cot(1/20)>799/40$ and $\tanh(5/2)>29/30$, so the arctangent argument is larger than $23171/1200$. Using $\arctan x>\pi/2-1/x$ and the elementary rational lower bound $\pi>333/106$, one gets
\[
 \int_{-50}^{50}K_B>\frac{967}{1000}.
\]
Thus
\begin{equation}\label{eq:highmid}
 R_{b,\ell}(\ell/2)
 >\frac{123}{125}\frac{967}{1000}
 =\frac{118941}{125000}>\frac{19}{20}.
\end{equation}
This proves \eqref{eq:HLH}.
\end{proof}

\section{All times}

\begin{proof}[Proof of Theorem~\ref{thm:main}]
Fix $t>0$ and define
\begin{equation}\label{eq:path}
 \ell(t)=\max\{100,95\pi t\},
 \qquad
 b(t)=\sqrt{\frac{\pi d t}{\ell(t)}},
 \qquad
 p(t)=\frac{8}{5\ell(t)},
 \qquad
 D(t)=\frac{d}{\ell(t)}.
\end{equation}
If $\ell(t)=95\pi t$, then $b(t)^2=d/95=1/100$. If $\ell(t)=100$, the defining inequality $95\pi t\le100$ again gives $b(t)^2\le1/100$. Therefore
\begin{equation}\label{eq:range}
 0<b(t)\le\frac1{10},\qquad \ell(t)\ge100,
 \qquad 0<p(t)\le\frac2{125}<1.
\end{equation}
Let $\nu_t=\nu_{b(t),p(t),\ell(t)}$ from \eqref{eq:witness}. Lemma~\ref{lem:crossings}, divided by $\ell(t)$, gives
\[
 S_t(0)>D(t),\qquad S_t(3)<D(t),\qquad S_t(\ell(t)/2)>D(t).
\]
Because $\ell(t)\ge100$, the three displayed points are strictly ordered. Also $S_t(y)\to0$ as $y\to\pm\infty$. Since $D(t)>0$ and $S_t$ is continuous, the equation
\begin{equation}\label{eq:Slevel}
 S_t(y)=D(t)
\end{equation}
has at least four distinct real roots: one to the left of $0$, one in $(0,3)$, one in $(3,\ell(t)/2)$, and one to the right of $\ell(t)/2$.

The parameter choice satisfies
\begin{equation}\label{eq:timeid}
 \frac{b(t)^2}{\pi D(t)}=t.
\end{equation}
Set
\[
 a(t)=\frac{D(t)}{b(t)}.
\]
Then
\[
 a(t)\pi t=b(t)<\pi,
\]
so $a(t)\in(0,1/t)$. For a root $y$ of \eqref{eq:Slevel}, put $r=e^{b(t)y}$. Since $q(by)=(A_b/b)S(y)$,
\[
 \frac{r}{c_b}\frac{d(\lambda_b\circledast\nu_t^{-1})}{dr}(r)
 =\frac{q(by)}{c_b}
 =\frac{A_bD(t)}{b(t)c_b}
 =\frac{\pi D(t)}{b(t)\sin b(t)}
 =\frac{a(t)\pi}{\sin(a(t)\pi t)}.
\]
Thus the four distinct roots of \eqref{eq:Slevel} yield four distinct positive roots of the source equation \eqref{eq:source-level}, contradicting the necessary condition for log-unimodality of $\sigma_t\boxtimes\nu_t$. Hence $\sigma_t\boxtimes\nu_t$ is not log-unimodal. Since $\nu_t$ is log-unimodal, $\sigma_t$ is not $\boxtimes$-strongly log-unimodal. As $t>0$ was arbitrary, the theorem follows.
\end{proof}

\section{Finite certificate and audit status}

The proof above is analytic; the accompanying script \texttt{check.py} uses exact rational arithmetic to certify the coarse numerical inequalities used in Lemma~\ref{lem:crossings}, including the rational $\pi$ enclosure, the three strict crossing margins, and the parameter identity $d/B^2=95$. Its archived output is \texttt{check.out}.

The source-matching theorem was independently audited end to end by two distinct repository tracks: R17 audit \texttt{R17-AUDIT-R18-I0080-C007-ENDTOEND-002} and P26 audit \texttt{P26-AUDIT-R18-I0080-C007-ENDTOEND-05}. Both audits returned full mathematical PASS verdicts and P26 independently re-executed the exact checker. These audits establish repository closure under the stated protocol; they are not substitutes for external peer review.

\begin{thebibliography}{9}
\bibitem{HUW}
T.~Hasebe, Y.~Ueda and J.-C.~Wang,
\emph{Log-unimodality for free positive multiplicative Brownian motion},
Colloq. Math. \textbf{169} (2022), 209--226;
arXiv:2009.13848 (2020).
\end{thebibliography}

\end{document}
